% Source: https://artofproblemsolving.com/community/c3h3639555_find_angle
% Stable identity: aops-topic-3639555-q01
Draw diagonal $AC$. Since $AD=DC$, $\angle DAC=\angle ACD=35^\circ$. Thus $\angle ACB=95^\circ$. By the sine rule in $ADC$, $AC=2\sin 55^\circ$ after scaling $AD=DC=CB=1$. In $ABC$, use the sine rule with $\angle ABC=85^\circ-\angle BAC$. It gives $2\sin 55^\circ\sin\angle BAC=\sin(85^\circ-\angle BAC)$. Expanding and solving, $\tan\angle BAC=\dfrac{\sin 85^\circ}{2\sin 55^\circ+\cos 85^\circ}=\dfrac{1}{\sqrt 3}$, using $2\sin 55^\circ+\cos 85^\circ=\sqrt 3\sin 85^\circ$. Since $0<\angle BAC<85^\circ$, it equals $30^\circ$, and $\boxed{\angle ABC=55^\circ}$.
\begin{asy}
import geometry; size(220);
pair D=(0,0),C=(1,0),A=dir(110),B=C+dir(50); draw(A--B--C--D--cycle); draw(A--C,dashed); label("$A$",A,NW);label("$B$",B,NE);label("$C$",C,SE);label("$D$",D,SW); label("$110^\circ$",D+(0.13,0.18));label("$130^\circ$",C+(-0.05,0.15));
\end{asy}
